\documentclass[11pt]{article}
\usepackage{amsmath,amssymb,amsthm}
\usepackage[margin=1in]{geometry}

\newtheorem{theorem}{Theorem}
\newtheorem{lemma}[theorem]{Lemma}
\newtheorem{corollary}[theorem]{Corollary}
\newtheorem{proposition}[theorem]{Proposition}
\theoremstyle{definition}
\newtheorem{definition}[theorem]{Definition}
\theoremstyle{remark}
\newtheorem{remark}[theorem]{Remark}

\newcommand{\orb}{\operatorname{orb}}
\newcommand{\supp}{\operatorname{supp}}
\newcommand{\Aalt}{A^{\mathrm{alt}}}
\newcommand{\thetaalt}{\theta^{\mathrm{alt}}}

\title{The Coset Poincar\'e Identity for the Top Atom Coefficient of $P_\mu(x;t)$}
\author{Clio}
\date{2026-07-28}

\begin{document}
\maketitle

\begin{abstract}
Let $\mu \vdash n$ and let $P_\mu(x;t)$ denote the Hall--Littlewood polynomial.
Expand $P_\mu$ in the basis of alternating Demazure atoms $\Aalt_\gamma$ ($\gamma$
ranging over the $S_n$-orbit of $\mu$):
$P_\mu = \sum_\gamma c_\gamma(t)\Aalt_\gamma$.
We prove
\[
  c_\mu(t) \;=\; \frac{[n]_t!}{\prod_i [m_i(\mu)]_t!}
             \;=\; \binom{n}{m_1(\mu),m_2(\mu),\ldots}_t
             \;=\; P(S_n/W_\mu)(t),
\]
where $W_\mu = \prod_i S_{m_i(\mu)}$ is the Young-subgroup stabilizer of $\mu$
and $\ell(\cdot)$ denotes Bruhat length on $\orb(\mu)$.
The proof rests on a \emph{second-moment invariant} of the support of atoms
together with an adjoint-pairing identity that isolates $c_\mu$ from all other
$c_\gamma$.
\end{abstract}

\section{Setup}

Fix a partition $\mu = (\mu_1 \ge \mu_2 \ge \cdots \ge \mu_n \ge 0) \vdash n$
in $n$ parts (with trailing zeros permitted).
For each $i \in \{0,1,2,\ldots\}$ let $m_i(\mu) = |\{j : \mu_j = i\}|$ be the
multiplicity of $i$ in $\mu$, and set $v_\mu(t) := \prod_i [m_i(\mu)]_t!$.

Let $\orb(\mu) := \{w \cdot \mu : w \in S_n\}$ be the $S_n$-orbit of $\mu$
(distinct rearrangements of $\mu$), acted on by
$(w \cdot \alpha)_j = \alpha_{w^{-1}(j)}$. For $\gamma \in \orb(\mu)$, the
Bruhat length $\ell(\gamma)$ is the minimum length of a sequence of adjacent
transpositions $s_{i_1},\ldots,s_{i_k}$ (bubble word) such that
$s_{i_k} \cdots s_{i_1} \cdot \mu = \gamma$; equivalently, the number of
inversions of $\gamma$ relative to $\mu$.

Define the alternating divided-difference operator
\[
  \thetaalt_i(f) \;:=\; \frac{x_{i+1} - t x_i}{x_i - x_{i+1}}\bigl(f - s_i f\bigr),
\]
where $s_i$ swaps $x_i \leftrightarrow x_{i+1}$. For a monomial $x^\alpha$:
if $\alpha_i = \alpha_{i+1}$ then $\thetaalt_i(x^\alpha) = 0$; if
$\alpha_i > \alpha_{i+1}$ (write $(a,b) = (\alpha_i,\alpha_{i+1})$) then
\begin{equation}\label{eq:theta-descent}
  \thetaalt_i(x^\alpha) \;=\; -t\,x^\alpha \;+\; (1-t)\sum_{b<j<a} x_i^{\,j} x_{i+1}^{\,a+b-j}\prod_{k \ne i,i+1} x_k^{\alpha_k} \;+\; x^{s_i\alpha}\,;
\end{equation}
if $\alpha_i < \alpha_{i+1}$ (write $(b,a) = (\alpha_i,\alpha_{i+1})$) then
\begin{equation}\label{eq:theta-ascent}
  \thetaalt_i(x^\alpha) \;=\; -x^\alpha \;-\; (1-t)\sum_{b<j<a} x_i^{\,j} x_{i+1}^{\,a+b-j}\prod_{k \ne i,i+1} x_k^{\alpha_k} \;+\; t\,x^{s_i\alpha}\,.
\end{equation}
The atom is defined by
\[
  \Aalt_\gamma \;:=\; \thetaalt_{i_k}\cdots\thetaalt_{i_1}(x^\mu),
\]
where $(i_1,\ldots,i_k)$ is a bubble word from $\mu$ to $\gamma$. In particular
$\Aalt_\mu = x^\mu$.

Finally, $P_\mu(x;t)$ satisfies $P_\mu = m_\mu + \sum_{\lambda \prec \mu} b_{\lambda,\mu}(t)\,m_\lambda$
in the monomial basis, where $\prec$ is strict dominance
(Macdonald III.\S 2).

We call an exponent $\alpha \in \mathbb Z_{\ge 0}^n$ an \emph{intermediate
of $x^\alpha$ at position $i$} if it agrees with $\alpha$ outside $(i,i+1)$ and
its $(i,i+1)$-entries are a pair $(j, a+b-j)$ with
$b < j < a$ where $\{a,b\} = \{\alpha_i,\alpha_{i+1}\}$, $a > b$.
Equations \eqref{eq:theta-descent}--\eqref{eq:theta-ascent} say the support of
$\thetaalt_i(x^\alpha)$ is $\{\alpha, s_i\alpha\}$ together with the intermediates
of $x^\alpha$ at $i$.

\section{Main theorem}

\begin{theorem}[Coset Poincar\'e identity]\label{thm:main}
For every partition $\mu \vdash n$,
\[
  c_\mu(t) \;=\; \frac{[n]_t!}{\prod_i [m_i(\mu)]_t!}
            \;=\; \binom{n}{m_1(\mu),m_2(\mu),\ldots}_t
            \;=\; P(S_n/W_\mu)(t).
\]
\end{theorem}

The theorem follows from three lemmas, proved in the next three subsections.

\subsection{The value $[x^\sigma] P_\mu$ on the orbit}

\begin{lemma}[Monomial-basis coefficient]\label{lem:mono}
For every $\sigma \in \orb(\mu)$, $[x^\sigma] P_\mu(x;t) = 1$.
\end{lemma}

\begin{proof}
Write $P_\mu = m_\mu + \sum_{\lambda \prec \mu} b_{\lambda,\mu}(t)\, m_\lambda$. Since
$m_\lambda = \sum_{\tau \in \orb(\lambda)} x^\tau$ and $\orb(\lambda) \cap \orb(\mu) = \emptyset$
for $\lambda \ne \mu$, we have $[x^\sigma] m_\lambda = 0$ for $\lambda \ne \mu$, while
$[x^\sigma] m_\mu = 1$ because $\sigma \in \orb(\mu)$.
\end{proof}

\subsection{The second-moment invariant}

Write $S_2(\alpha) := \sum_j \alpha_j^2$.

\begin{lemma}[Second-moment invariant]\label{lem:secondmoment}
For every $\gamma \in \orb(\mu)$ and every $\alpha \in \supp(\Aalt_\gamma)$,
\[
  S_2(\alpha) \;\le\; S_2(\mu).
\]
Moreover, $S_2$ is strictly decreased by every ``intermediate'' contribution
in \eqref{eq:theta-descent}--\eqref{eq:theta-ascent} and left unchanged by the
$\alpha$-diagonal and $s_i\alpha$ contributions.
\end{lemma}

\begin{proof}
By induction on $\ell(\gamma)$.

\emph{Base.} $\gamma = \mu$: $\supp(\Aalt_\mu) = \{\mu\}$ and $S_2(\mu) \le S_2(\mu)$. \checkmark

\emph{Step.} $\gamma = s_i \gamma'$ with $\ell(\gamma') = \ell(\gamma)-1$, so
$\Aalt_\gamma = \thetaalt_i(\Aalt_{\gamma'})$. Fix $\alpha \in \supp(\Aalt_\gamma)$;
then $\alpha \in \supp(\thetaalt_i(x^\beta))$ for some $\beta \in \supp(\Aalt_{\gamma'})$
with $\beta_i \ne \beta_{i+1}$.

Case (i): $\alpha = \beta$. Then $S_2(\alpha) = S_2(\beta) \le S_2(\mu)$ by IH.

Case (ii): $\alpha = s_i\beta$. Same multiset as $\beta$, so
$S_2(\alpha) = S_2(\beta) \le S_2(\mu)$.

Case (iii): $\alpha$ is an intermediate of $x^\beta$ at $i$. Write
$\{a,b\} = \{\beta_i,\beta_{i+1}\}$ with $a>b$, and $(\alpha_i,\alpha_{i+1}) = (j,a+b-j)$
(or the reverse) with $b<j<a$. Then
\begin{align*}
  S_2(\alpha) - S_2(\beta)
  &= j^2 + (a+b-j)^2 - a^2 - b^2 \\
  &= -\bigl[(a^2+b^2) - (j^2 + (a+b-j)^2)\bigr] \\
  &= -\,2(a-j)(j-b) \;<\; 0
\end{align*}
because $b < j < a$. So $S_2(\alpha) < S_2(\beta) \le S_2(\mu)$.
\qedhere
\end{proof}

The last identity---$(a^2+b^2) - (j^2 + (a+b-j)^2) = 2(a-j)(j-b)$---is the
crucial arithmetic fact. It says: replacing an outer pair $\{a,b\}$ by an
inner pair $\{j, a+b-j\}$ with the same sum \emph{strictly reduces} the sum
of squares.

\subsection{The Weyl-symmetrization identity}

Define the linear functional $L: \mathbb Z[t][x_1,\ldots,x_n] \to \mathbb Z[t]$ by
\[
  L(f) \;:=\; \sum_{\sigma \in \orb(\mu)} t^{\ell(\sigma)}\, [x^\sigma]\, f.
\]
Equivalently, $L(f) = \langle f, Q\rangle$ where $Q := \sum_{\sigma \in \orb(\mu)} t^{\ell(\sigma)} x^\sigma$
and $\langle x^\alpha, x^\beta\rangle := \delta_{\alpha,\beta}$ is the standard
monomial pairing.

\begin{lemma}[Weyl-symmetrization identity]\label{lem:weyl}
For every $\gamma \in \orb(\mu)$,
\[
  L(\Aalt_\gamma) \;=\; \delta_{\gamma,\mu}.
\]
\end{lemma}

\begin{proof}
Induction on $\ell(\gamma)$.

\emph{Base.} $\gamma = \mu$: $\Aalt_\mu = x^\mu$, so
$L(\Aalt_\mu) = t^{\ell(\mu)} = t^0 = 1$. \checkmark

\emph{Step.} $\gamma = s_i \gamma'$ with $\ell(\gamma') < \ell(\gamma)$. We show
$L(\thetaalt_i(\Aalt_{\gamma'})) = 0$.

Let $\widetilde\theta_i$ denote the adjoint of $\thetaalt_i$ with respect to $\langle,\rangle$:
$\langle \thetaalt_i f, g\rangle = \langle f, \widetilde\theta_i g\rangle$, i.e.
$[x^\alpha]\widetilde\theta_i(x^\beta) = [x^\beta]\thetaalt_i(x^\alpha)$. Then
\[
  L(\thetaalt_i(\Aalt_{\gamma'})) \;=\; \langle \thetaalt_i(\Aalt_{\gamma'}), Q\rangle
                                 \;=\; \langle \Aalt_{\gamma'}, \widetilde\theta_i(Q)\rangle.
\]
We evaluate $\widetilde\theta_i(Q)$ and pair it against $\Aalt_{\gamma'}$.

\medskip\noindent\emph{Step (a): the orbit-supported part of $\widetilde\theta_i(Q)$ vanishes.}

Reading off \eqref{eq:theta-descent}--\eqref{eq:theta-ascent}, for each
$\sigma \in \orb(\mu)$ the ``diagonal-and-swap'' contribution to $\widetilde\theta_i(x^\sigma)$ is
\[
  \begin{cases}
    -t\,x^\sigma + t\,x^{s_i\sigma} & \text{if $\sigma_i > \sigma_{i+1}$ (descent)},\\
    -x^\sigma + x^{s_i\sigma} & \text{if $\sigma_i < \sigma_{i+1}$ (ascent)},\\
    0 & \text{if $\sigma_i = \sigma_{i+1}$}.
  \end{cases}
\]
Indeed, $[x^\sigma]\thetaalt_i(x^\sigma)$ contributes at $\alpha = \sigma$
(giving $-t$ or $-1$ or $0$), and $[x^\sigma]\thetaalt_i(x^{s_i\sigma})$ contributes
at $\alpha = s_i\sigma$ (giving $t$ or $1$).

Descent-ascent pairs $(\sigma, s_i\sigma)$ in the orbit satisfy
$\ell(s_i\sigma) = \ell(\sigma)+1$ when $\sigma$ has a descent at $(i,i+1)$
(swapping a descent up the Bruhat order). Grouping $\sigma$ (descent) with
$\sigma' := s_i\sigma$ (ascent), the diagonal-and-swap contributions to
$\widetilde\theta_i(Q) = \sum_\sigma t^{\ell(\sigma)} \widetilde\theta_i(x^\sigma)$ at $x^\sigma$ and $x^{\sigma'}$
are:
\begin{align*}
  {[x^\sigma]}: \quad & t^{\ell(\sigma)}(-t) + t^{\ell(\sigma')}(1) \;=\; -t^{\ell(\sigma)+1} + t^{\ell(\sigma)+1} \;=\; 0,\\
  {[x^{\sigma'}]}: \quad & t^{\ell(\sigma)}(t) + t^{\ell(\sigma')}(-1) \;=\; t^{\ell(\sigma)+1} - t^{\ell(\sigma)+1} \;=\; 0.
\end{align*}
Orbit elements fixed by $s_i$ contribute $0$ directly.
Hence the orbit-supported part of $\widetilde\theta_i(Q)$ vanishes.

\medskip\noindent\emph{Step (b): the non-orbit part of $\widetilde\theta_i(Q)$ has strictly
higher second moment.}

The remaining contribution to $\widetilde\theta_i(x^\sigma)$, for each $\sigma \in \orb(\mu)$,
comes from ``reverse intermediates'': $\alpha \ne \sigma, s_i\sigma$ such that $\sigma$
is an intermediate of $\thetaalt_i(x^\alpha)$.
Reading off \eqref{eq:theta-descent}--\eqref{eq:theta-ascent}, such $\alpha$ has
$(\alpha_i,\alpha_{i+1}) = (u,v)$ or $(v,u)$ with $u > v$, $u + v = \sigma_i + \sigma_{i+1}$,
$v < \min(\sigma_i,\sigma_{i+1})$, $u > \max(\sigma_i,\sigma_{i+1})$, and
$\alpha_k = \sigma_k$ for $k \ne i,i+1$. Set
$M := \max(\sigma_i,\sigma_{i+1})$ and $M' := \min(\sigma_i,\sigma_{i+1})$
(so $M + M' = u + v$), and write $p := u - M \ge 1$ (integer). Then $v = M' - p$,
and the constraint $v \ge 0$ says $p \le M'$. Consequently,
\begin{align*}
  S_2(\alpha) - S_2(\sigma)
  &= (u^2 + v^2) - (M^2 + {M'}^2)\\
  &= (M+p)^2 + (M'-p)^2 - M^2 - {M'}^2\\
  &= 2p(M - M') + 2p^2 \;\ge\; 2p^2 \;\ge\; 2 \;>\; 0
\end{align*}
(using $M \ge M'$ and $p \ge 1$). Since $\sigma \in \orb(\mu)$ implies
$S_2(\sigma) = S_2(\mu)$, every reverse-intermediate $\alpha$ satisfies
\begin{equation}\label{eq:strict-inequality}
  S_2(\alpha) \;>\; S_2(\mu).
\end{equation}

\medskip\noindent\emph{Step (c): the pairing vanishes.}

By Lemma~\ref{lem:secondmoment}, every $\alpha \in \supp(\Aalt_{\gamma'})$
satisfies $S_2(\alpha) \le S_2(\mu)$. In particular
$[x^\alpha]\Aalt_{\gamma'} = 0$ for every $\alpha$ with $S_2(\alpha) > S_2(\mu)$.
By step (a), the orbit-supported part of $\widetilde\theta_i(Q)$ vanishes; and by
step (b) plus \eqref{eq:strict-inequality}, every non-orbit monomial in the
support of $\widetilde\theta_i(Q)$ has $S_2(\alpha) > S_2(\mu)$, hence pairs to zero
against $\Aalt_{\gamma'}$. Therefore
\[
  L(\Aalt_\gamma) \;=\; \langle \Aalt_{\gamma'}, \widetilde\theta_i(Q) \rangle \;=\; 0,
\]
completing the induction.
\end{proof}

\subsection{Proof of Theorem~\ref{thm:main}}

Applying $L$ to $P_\mu = \sum_{\gamma \in \orb(\mu)} c_\gamma(t)\,\Aalt_\gamma$ and using
Lemma~\ref{lem:weyl}:
\[
  L(P_\mu) \;=\; \sum_\gamma c_\gamma(t)\, L(\Aalt_\gamma) \;=\; c_\mu(t).
\]
On the other hand, using Lemma~\ref{lem:mono},
\[
  L(P_\mu) \;=\; \sum_{\sigma \in \orb(\mu)} t^{\ell(\sigma)}\,[x^\sigma] P_\mu
             \;=\; \sum_{\sigma \in \orb(\mu)} t^{\ell(\sigma)}
             \;=\; P(S_n/W_\mu)(t).
\]
The last equality is standard: the parabolic Poincar\'e polynomial
$P(S_n/W_\mu)(t) = \sum_{w \in D_\mu} t^{\ell(w)}$, where $D_\mu$ is the set of
minimum-length coset representatives for $S_n/W_\mu$; the map $w \mapsto w \cdot \mu$
is a bijection $D_\mu \to \orb(\mu)$ preserving length. Finally,
\[
  P(S_n/W_\mu)(t) \;=\; \binom{n}{m_1(\mu),m_2(\mu),\ldots}_t
                  \;=\; \frac{[n]_t!}{\prod_i [m_i(\mu)]_t!}.
\]
Combining, $c_\mu(t) = [n]_t!/v_\mu(t)$. \qed

\section{Verification}

\subsection{Computational evidence}

Verified numerically (via SageMath and SymPy scripts) on:
\begin{itemize}\itemsep=0pt
  \item Every partition of $n$ with $n \le 5$ that has been tested previously
        (12 partitions: $(2,1,0), (2,2,0), (2,2,1), (3,2,1), (2,2,0,0), (2,2,1,1), (3,3,1,0), (3,3,2,0), (3,2,2,0), (3,2,2,1), (2,2,2,1,0), (2,2,2,1,1)$).
  \item Additional test cases including $(3,2,1,0), (4,2,1,0), (4,3,2,1,0), (4,2,2,0)$
        checking that (i) the atom-triangularity lemma holds, (ii) the second-moment
        invariant holds (equality attained only at $\alpha = \sigma \in \orb(\mu)$),
        (iii) the Weyl-symmetrization identity $L(\Aalt_\gamma) = \delta_{\gamma,\mu}$ holds,
        (iv) the resulting $c_\mu$ matches $[n]_t!/v_\mu(t)$.
\end{itemize}
Zero failures across all cases.

\subsection{Worked example: $\mu = (2,2,0)$}

Orbit is $\{(2,2,0), (2,0,2), (0,2,2)\}$ with Bruhat lengths $0,1,2$.

\emph{Atoms.} By direct computation via \eqref{eq:theta-descent}:
\begin{align*}
  \Aalt_{(2,2,0)} &= x_1^2 x_2^2,\\
  \Aalt_{(2,0,2)} &= -t\,x_1^2 x_2^2 + (1-t)\,x_1^2 x_2 x_3 + x_1^2 x_3^2,\\
  \Aalt_{(0,2,2)} &= -t(1-t)\,x_1^2 x_2 x_3 + (1-t)\,x_1 x_2^2 x_3 - t\,x_1^2 x_3^2 + (1-t)\,x_1 x_2 x_3^2 + x_2^2 x_3^2.
\end{align*}
Second moments:
\begin{itemize}\itemsep=0pt
\item $\Aalt_{(2,2,0)}$: only $S_2 = 8 = S_2(\mu)$.
\item $\Aalt_{(2,0,2)}$: monomials $(2,2,0), (2,1,1), (2,0,2)$ have $S_2 = 8, 6, 8$; all $\le 8$.
\item $\Aalt_{(0,2,2)}$: monomials $(2,1,1),(1,2,1),(2,0,2),(1,1,2),(0,2,2)$ have $S_2 = 6,6,8,6,8$; all $\le 8$.
\end{itemize}

\emph{The identity $L(\Aalt_\gamma) = \delta_{\gamma,\mu}$:}
\begin{align*}
  L(\Aalt_{(2,2,0)}) &= t^0 \cdot 1 = 1.\\
  L(\Aalt_{(2,0,2)}) &= t^0 (-t) + t^1 (0) + t^2 (0) + \text{(non-orbit terms $\to$ 0)} \\
                     &\quad \phantom{=} \text{[explicitly: only } (2,2,0)\text{ and } (2,0,2) \text{ are in orbit:]}\\
                     &= 1 \cdot (-t) + t \cdot 1 = 0.\\
  L(\Aalt_{(0,2,2)}) &= t^0 \cdot 0 + t^1 (-t) + t^2 \cdot 1 = -t^2 + t^2 = 0.
\end{align*}

\emph{Verification of $c_\mu$:} solving $c_\gamma\cdot M = (1,1,1)$ by back-substitution
(with $M$ upper triangular, $M[\gamma,\sigma] = [x^\sigma]\Aalt_\gamma$, ordered
by decreasing $\ell$):
$c_{(0,2,2)}=1$, $c_{(2,0,2)}=1+t=[2]_t$, $c_{(2,2,0)}=1+t+t^2=[3]_t$.
And indeed $c_\mu = [3]_t = [3]_t!/(v_\mu = [2]_t) = [3]_t!\,/\,[2]_t \cdot [1]_t = [3]_t$. \checkmark

\section{Remarks}

\begin{remark}
The proof uses only the following inputs:
\begin{enumerate}\itemsep=0pt
\item The definition of $\thetaalt_i$ and the associated atom construction.
\item The monomial expansion $P_\mu = m_\mu + $ strictly-lower-dominance terms.
\item The parabolic Poincar\'e polynomial identity
      $\sum_{w \in D_\mu} t^{\ell(w)} = [n]_t!/v_\mu(t)$ (Chevalley--Shephard--Todd).
\end{enumerate}
The two novel structural facts---the second-moment invariant
(Lemma~\ref{lem:secondmoment}) and the Weyl-symmetrization identity
(Lemma~\ref{lem:weyl})---may be of independent interest.
\end{remark}

\begin{remark}
Lemma~\ref{lem:weyl} is strictly stronger than the top-coefficient identity:
it gives a \emph{characterization} of the top atom by a linear functional. In
particular, $L$ picks out $c_\mu$ exactly, while all other $c_\gamma$ are
annihilated. This suggests $L$ may extend to a duality between the atom basis
and a suitably-graded orbit basis, which we intend to explore.
\end{remark}

\begin{remark}
The second-moment invariant reflects a general geometric fact: the operation
of ``pulling in'' an outer pair $(a,b)$ with $a>b$ to an inner pair $(j, a+b-j)$
with $b<j<a$ is strictly a convex-combination move (it decreases the ``spread''
of the multiset while preserving the mean). The formal identity
$(a^2+b^2) - (j^2 + (a+b-j)^2) = 2(a-j)(j-b)$
is thus a manifestation of convexity. This ties the atom-support structure to
a classical inequality.
\end{remark}

\end{document}
